%% MSR 2027 paper draft — Issue-triage Agent Skills in GitHub Agentic Workflows
%% Overleaf: upload this folder; set main document to main.tex
%% Requires ACM `acmart` class (available on Overleaf).

\documentclass[sigconf,review,anonymous]{acmart}

\usepackage{xurl}
\makeatletter
\g@addto@macro\UrlBreaks{\do\.\do\-\do\:\do\*\do\~}
\makeatother
\newcommand{\filepath}[1]{{\path{#1}}}

\AtBeginDocument{%
  \providecommand\BibTeX{{Bib\TeX}}}

\setcopyright{acmlicensed}
\copyrightyear{2026}
\acmYear{2026}
\acmDOI{XXXXXXX.XXXXXXX}
\acmConference[MSR '27]{Mining Software Repositories}{August 2026}{}{}
\acmISBN{978-1-4503-XXXX-X/26/08}

\title{What Makes Issue-Triage Skills Effective in GitHub Agentic Workflows?}

\author{Dung Pham(s)}
\affiliation{%
  \institution{Trent University}
  \city{Peterborough}
  \country{Canada}}
\email{dungpham@trentu.ca}

\renewcommand{\shortauthors}{Anonymous et al.}

\begin{abstract}
Issue triage is among the most common tasks in public GitHub Agentic Workflows
(\emph{gh-aw}): agents classify issues, apply or suggest labels, detect duplicates,
and route work to maintainers.
Despite this prevalence, few triage workflows attach reusable Agent Skills
(\texttt{SKILL.md}, scripts, and references); most encode policy only as inline
prompts.
We mine gh-aw workflow markdown, skill packages, and run artifacts to characterize
\emph{what makes a triage skill good}---procedural specificity, label taxonomies,
duplicate and spam rules, approval constraints, and sibling artifacts---and we link
those properties to observable GitHub outcomes such as label accuracy and maintainer
override rate.
Building on those findings, we propose an \textbf{issue-triage skill generator}
(or equivalent custom instructions) that produces repository-specific skill packages.
Generated skills are meant to be consumed by \texttt{gh aw compile}, so the resulting
\filepath{.lock.yml} installs and activates them in the compiled agentic pipeline.
\end{abstract}

\begin{CCSXML}
<ccs2012>
 <concept>
  <concept_id>10011007.10011006.10011008</concept_id>
  <concept_desc>Software and its engineering~Software maintenance</concept_desc>
  <concept_significance>500</concept_significance>
 </concept>
 <concept>
  <concept_id>10010147.10010257</concept_id>
  <concept_desc>Computing methodologies~Machine learning</concept_desc>
  <concept_significance>300</concept_significance>
 </concept>
</ccs2012>
\end{CCSXML}

\ccsdesc[500]{Software and its engineering~Software maintenance}
\ccsdesc[300]{Computing methodologies~Machine learning}

\keywords{Agent Skills, Issue Triage, GitHub Agentic Workflows, Mining Software Repositories, LLM Agents}

\begin{document}

\maketitle

%==============================================================================
\section{Introduction}
\label{sec:introduction}

Coding agents increasingly rely on \emph{Agent Skills}---packaged procedural knowledge
(\texttt{SKILL.md}), scripts, and reference files---to perform recurring software tasks.
GitHub Agentic Workflows (\emph{gh-aw}) compile markdown workflow definitions into
executable CI pipelines (\filepath{.lock.yml}) that run agents inside GitHub Actions.

Among mined public gh-aw deployments, \textbf{issue triage} is a dominant task:
agents read an existing issue, classify it against a repository-specific taxonomy,
apply or suggest labels, optionally detect duplicates or spam, and post comments or
update project boards.
This is not a trivial labeling rule.
Effective triage requires duplicate judgment, conservative spam handling, allowlisted
labels, and---in some repositories---human approval via GitHub issue intents rather
than immediate writes.

\textbf{Problem.}
Despite the complexity of triage, most public triage workflows do \emph{not} attach
skills.
They encode policy as long inline prompts.
Where skills exist (e.g., \texttt{cli/cli}, \texttt{mono/SkiaSharp}), they externalize
taxonomies, duplicate-detection procedures, and safety constraints that used to live
in maintainer handbooks.
Practitioners therefore lack evidence-based guidance on \emph{what makes a triage
skill good}, and they lack tooling to \emph{author} such a skill and compile it into
a working gh-aw pipeline.

\textbf{Main question.}
\emph{What characteristics of Agent Skills make GitHub Agentic Workflows effective
at issue triage---and can those characteristics be generated and compiled?}
This breaks down into:
\begin{enumerate}
  \item What does \emph{good} mean for an issue-triage skill in gh-aw?
  \item Which intrinsic and operational properties predict triage quality on GitHub
        (labels, comments, approval, consistency)?
  \item Can we generate repository-specific triage skills (or custom instructions)
        from mined patterns, and can \texttt{gh aw compile} include them in the
        resulting \filepath{.lock.yml}?
\end{enumerate}

Prior skill benchmarks emphasize task success on isolated tests~\cite{utilitybench2026,skillsbench2026}.
SkillEval~\cite{skilleval2026} advances \emph{intrinsic} evaluation of
\texttt{SKILL.md} text.
We differentiate by (i)~restricting the task to \textbf{issue triage in gh-aw},
(ii)~evaluating not only \texttt{SKILL.md} but also sibling artifacts and
\emph{how} skills are registered, installed, invoked, and followed in run
artifacts, and (iii)~closing the loop with a \textbf{generator} whose output is a
first-class input to the gh-aw compiler.

\textbf{Contributions.}
\begin{itemize}
  \item A mining study of public gh-aw issue-triage workflows, including how often
        they attach skills, how they encode taxonomy and approval, and how skills
        are actually picked up at runtime.
  \item An evaluation framework for triage skill quality: intrinsic features,
        artifact siblings, five-stage skill pickup, and GitHub outcome metrics
        (label accuracy, override rate, time-to-triage, false spam/duplicate).
  \item Research questions linking those properties to triage effectiveness.
  \item A proposed \textbf{issue-triage skill generator} (or custom-instruction
        pack) whose generated \texttt{SKILL.md} (and references) are declared in
        workflow frontmatter so that \texttt{gh aw compile} emits a
        \filepath{.lock.yml} that installs and activates the skill.
\end{itemize}

%==============================================================================
\section{Motivation and Scope}
\label{sec:scope}

\paragraph{Focus scope.}
We scope the study to \textbf{issue-triage agentic workflows} compiled by gh-aw:
classifying issues, labeling (or suggesting labels), duplicate and spam handling,
comments, project-board updates, and related approval gates.
We study skills that those workflows attach---via frontmatter \texttt{skills:},
inline \texttt{\#\#\ skill:} blocks, hint/fusion fragments, or prompt-side paths
such as \filepath{.github/skills/} and \filepath{.agents/skills/}---rather than
skills in general coding-agent corpora.

CI investigation and fix-and-PR workflows appear in the broader gh-aw ecosystem
and may share skill-authoring lessons, but they are out of scope for the generator
and for the primary outcome metrics.

\paragraph{Why triage, and why skills.}
Before gh-aw, repositories handled triage with issue templates, keyword bots,
Probot apps, and maintainer meetings.
Those mechanisms are brittle on duplicates, taxonomy nuance, and spam.
Agentic triage can reason over issue text, but without a skill the policy lives
only in the workflow markdown and is hard to reuse, test, or compile consistently.
A good skill should encode the maintainer handbook: allowed labels, when to
suggest versus apply, how to search for duplicates, and when to stop.

\paragraph{From findings to a compiler-ready generator.}
Our planned methodology has four stages:
\begin{enumerate}
  \item \textbf{Mine} public gh-aw triage workflows and attached skill packages.
  \item \textbf{Characterize} properties that distinguish effective triage skills
        (specificity, taxonomy files, failure handling, modularity, approval
        constraints) and correlate them with GitHub outcomes.
  \item \textbf{Generate} a repository-specific skill (or custom instructions)
        from those patterns plus repo context (existing labels, issue templates,
        safe-output allowlists).
  \item \textbf{Compile} the generated skill through \texttt{gh aw compile} so
        the lock file contains activation steps (\texttt{gh skill install},
        checkout of skill paths, engine catalog) rather than leaving skills as
        undocumented prompt text.
\end{enumerate}

%==============================================================================
\section{Background and Related Work}
\label{sec:related}

\paragraph{Issue triage before agentic workflows.}
Classical OSS triage used templates and forms, default \texttt{needs-triage}
labels, GitHub Actions keyword rules, stale bots, and GitHub Apps.
Human maintainers remained the ground truth for duplicates and routing.
Agentic workflows add a middle layer: an LLM agent with bounded
\texttt{safe-outputs} (e.g., \texttt{add-labels}, \texttt{add-comment}) rather
than unrestricted \texttt{GITHUB\_TOKEN} writes.

\paragraph{SkillsBench and task-based evaluation.}
SkillsBench~\cite{skillsbench2026} evaluates agent skills at scale with emphasis on
task completion.
Related work on skill selection, retrieval, and adaptation~\cite{adaptation2026}
studies how agents choose and refine skills during execution.
Utility-oriented benchmarks~\cite{utilitybench2026} primarily report task success
rate.
We instead measure \textbf{triage effectiveness on GitHub}: label agreement with
maintainer ground truth, suggestion acceptance, and workflow-run conclusions.

\paragraph{SkillEval (intrinsic evaluation).}
SkillEval~\cite{skilleval2026} compares skills along four dimensions with seven
intrinsic metrics.
We extend this line of work by analyzing sibling artifacts (scripts, references)
and by tying scores to \emph{triage} outcomes rather than generic task success.

\paragraph{GitHub Agentic Workflows and skills.}\mbox{}\par\noindent
gh-aw compiles \filepath{.github/workflows/*.md} into lock files and runs agents
inside GitHub Actions~\cite{ghaw2026}.
Skills enter a workflow through compiler-supported mechanisms~\cite{ghawFrontmatter2026,ghawSkillsGuide2026,ghawFaq2026}:
frontmatter \texttt{skills:} (preferred), inline \texttt{\#\#\ skill:} blocks,
prompt-side \emph{hint} and \emph{fusion} fragments, and (separately)
\texttt{plugins:} for Agent Plugins.
Copilot discovers on-disk skills from
\filepath{.github/skills/}, \filepath{.claude/skills/}, and
\filepath{.agents/skills/} and loads \texttt{SKILL.md} on demand~\cite{copilotAboutSkills2026,copilotAddSkillsCli2026}.
Each skill follows the open \texttt{SKILL.md} specification~\cite{agentskillsSpec2026}.
Workflow runs produce artifacts that distinguish skill \emph{declaration},
\emph{installation}, \emph{availability}, and \emph{use}.

\paragraph{GitSkills dataset.}
The GitSkills corpus~\cite{gitskills2026} catalogs skill artifacts across
repositories---our source for intrinsic feature extraction, joined to gh-aw
triage workflows that actually reference those packages.

%==============================================================================
\section{Research Questions}
\label{sec:rqs}

\begin{description}
  \item[RQ1] How are issue-triage gh-aw workflows structured, and how often do
        they attach Agent Skills versus inline prompts only?
  \item[RQ2] What characteristics distinguish high-quality issue-triage skills
        (and skill-using workflows) from low-quality or skill-free ones?
  \item[RQ3] Which intrinsic properties---length, procedural specificity, taxonomy
        completeness, failure handling, modularity of scripts/references---predict
        triage outcomes on GitHub?
  \item[RQ4] How do practitioners gate triage actions (auto-apply labels versus
        \texttt{issue-intent} suggestions versus second-agent comment review),
        and how should a generated skill encode those constraints?
  \item[RQ5] Can mined patterns plus repository context generate an issue-triage
        skill (or custom instructions) that improves labeling accuracy and
        consistency relative to inline-prompt baselines?
  \item[RQ6] Does declaring the generated skill in workflow frontmatter and running
        \texttt{gh aw compile} produce a \filepath{.lock.yml} that correctly
        installs, catalogs, and enables the skill at activation time?
\end{description}

RQ6 is the compiler-facing success criterion: generation is incomplete unless the
skill survives compilation into a runnable lock file.

%==============================================================================
\section{Methodology}
\label{sec:methodology}

\subsection{Identifying gh-aw repositories and triage workflows}
We detect gh-aw usage when a public repository contains a matching workflow pair
on the default branch:
\filepath{.github/workflows/\{name\}.md} and
\filepath{.github/workflows/\{name\}.lock.yml}.
We crawl candidates via the GitHub Code Search API with size-based query
splitting to overcome the 1{,}000-result cap.

We classify a workflow as \textbf{issue triage} from workflow name, frontmatter
\texttt{description}, and prompt body (keywords such as triage, labeler, duplicate
detection, issue classification), then extract unique repositories and
repository--workflow pairs across crawl snapshots.

\subsection{How gh-aw registers skills}
\label{sec:skill-registration}

gh-aw separates \textbf{skill installation} (compiler-managed, before the agent
job) from \textbf{prompt shaping} (how skill text reaches the model).
We treat these as distinct signals.

\paragraph{Frontmatter \texttt{skills:} (preferred).}\mbox{}\par\noindent
The top-level \texttt{skills:} array instructs gh-aw to install Copilot skills
during the \emph{activation} job~\cite{ghawFrontmatter2026,ghawSkillsGuide2026}.
Entries may be local paths (e.g., \filepath{.github/skills/issue-triage}) or
external \texttt{owner/repo[/path]@ref} references.
The compiler emits activation steps; authors must not add manual
\texttt{gh skill install} steps~\cite{ghawSkillsGuide2026}.
This is the path our generator must produce so \texttt{gh aw compile} includes
the skill in \filepath{.lock.yml}.

\paragraph{Inline, hint, and fusion strategies.}
Beyond \texttt{skills:}, gh-aw supports prompt-side skill content:
\begin{itemize}
  \item \textbf{Inline skills} --- \texttt{\#\#\ skill: `name`} blocks extracted
        at setup into engine skill directories~\cite{ghawSkillsGuide2026};
  \item \textbf{Hint} --- prompt text directing the agent to discover
        \texttt{SKILL.md} under \filepath{skills/} or \filepath{.github/skills/};
  \item \textbf{Fusion} --- comments that splice \texttt{SKILL.md} sections into
        the workflow prompt~\cite{ghawSkillsGuide2026}.
\end{itemize}
These shape the prompt but do not replace \texttt{skills:} installation for
packages the compiler should pin.

\paragraph{Plugins and APM imports.}
\texttt{plugins:} and APM \texttt{imports:}/\texttt{packages:} can distribute
skills~\cite{ghawFaq2026}.
We record them separately; primary labels target \texttt{SKILL.md}-based Copilot
skills used for triage.

\subsection{How Copilot discovers skills at runtime}
\label{sec:skill-discovery}

After checkout and installation, Copilot CLI discovers project skills from
\filepath{.github/skills/}, \filepath{.claude/skills/}, or
\filepath{.agents/skills/}~\cite{copilotAboutSkills2026,copilotAddSkillsCli2026}.
The process log exposes \texttt{<available\_skills>} with
\texttt{<location>project|builtin</location>}.
Availability in the catalog does not imply the skill body was loaded or
followed~\cite{copilotAboutSkills2026}.

The task prompt is the markdown body of the workflow file; the final agent prompt
is \filepath{aw-prompts/prompt.txt}.
Frontmatter \texttt{skills:} are not copied into that file automatically; they
reach the model through discovery plus optional hint/fusion/inline text.

\subsection{Operationalizing skill pickup from run artifacts}
\label{sec:skill-extraction}

We operationalize five non-equivalent stages (Table~\ref{tab:skill-stages}).

\begin{table}[t]
  \caption{Skill lifecycle stages and primary evidence in gh-aw run artifacts.}
  \label{tab:skill-stages}
  \small
  \begin{tabular}{@{}p{0.14\linewidth}p{0.30\linewidth}p{0.48\linewidth}@{}}
    \toprule
    \textbf{Stage} & \textbf{Question} & \textbf{Primary artifact / signal} \\
    \midrule
    Configured &
    Declared in workflow? &
    \filepath{.github/workflows/*.md} frontmatter \texttt{skills:}, inline
    \texttt{\#\#\ skill:}, hint/fusion; \filepath{activation/aw_info.json}
    \texttt{"skills"} list \\
    Installed &
    On disk before agent? &
    Activation console (\texttt{gh skill install}); \texttt{SKILL.md} under
    \filepath{activation/}; \filepath{pre-agent-audit.txt} \\
    Available &
    In Copilot catalog? &
    \filepath{process-*.log} $\rightarrow$ \texttt{<available\_skills>} with
    \texttt{project} location \\
    Invoked &
    Skill tool called? &
    Same log $\rightarrow$ \texttt{tool\_calls} with \texttt{"name": "skill"} \\
    Followed &
    Instructions in output? &
    Fingerprints from \texttt{SKILL.md} in agent output, labels, comments \\
    \bottomrule
  \end{tabular}
\end{table}

For triage, \textbf{Followed} includes whether applied or suggested labels match
the skill's taxonomy, whether duplicate comments cite issues as instructed, and
whether \texttt{issue-intent} suggestions (versus direct apply) match the skill's
approval policy.

\subsection{Feature extraction for triage skills}
From attached skill packages and workflow prompts we extract:
\begin{itemize}
  \item \texttt{SKILL.md} length, frontmatter validity, procedural steps;
  \item presence of a \textbf{label taxonomy} (allowlist, mutually exclusive
        types, area labels);
  \item duplicate-search and spam/safety procedures;
  \item approval encoding (\texttt{issue-intent: true}, suggest-only language,
        dry-run, second-agent comment review);
  \item modularity (scripts, \filepath{references/}, sibling count);
  \item safe-output bounds (\texttt{add-labels} max, allowed labels).
\end{itemize}

\subsection{Outcome metrics for issue triage}
Where GitHub data permit, we measure:
\begin{itemize}
  \item \textbf{Label accuracy} --- agreement with maintainer-final labels;
  \item \textbf{Override / rejection rate} --- suggested intents not accepted;
  \item \textbf{Time-to-triage} --- open to \texttt{needs-triage} removal or
        \texttt{triage/triaged};
  \item \textbf{False spam / false duplicate} --- reversed by humans;
  \item \textbf{Workflow conclusion} --- gh-aw run success and safe-output
        emission;
  \item \textbf{Pickup completeness} --- configured through followed
        (Table~\ref{tab:skill-stages}).
\end{itemize}

\subsection{Issue-triage skill generator}
\label{sec:generator}

The generator is the constructive contribution of the study.
It takes as input:
\begin{itemize}
  \item mined \emph{patterns} from high-quality triage skills (structure of
        \texttt{SKILL.md}, taxonomy files, duplicate/spam sections, approval
        language);
  \item \emph{repository context}: existing labels, issue templates, CODEOWNERS,
        current workflow \texttt{safe-outputs} allowlists.
\end{itemize}
It emits:
\begin{itemize}
  \item a skill directory (e.g., \filepath{.github/skills/issue-triage/SKILL.md}
        plus \filepath{references/label-taxonomy.md});
  \item optional custom-instruction fragments for the workflow prompt body
        (hints to load the skill, never apply labels outside the allowlist);
  \item a patch to workflow frontmatter adding
        \texttt{skills: [.github/skills/issue-triage]} (or the generated path).
\end{itemize}

\paragraph{Compile integration.}
\texttt{gh aw compile} must see a valid \texttt{skills:} (or inline skill)
declaration so the lock file contains activation jobs that install the generated
package, sparse-checkout skill paths if needed, and leave \texttt{SKILL.md} on
disk for Copilot discovery~\cite{ghawSkillsGuide2026}.
We treat a compile as successful for RQ6 when:
(1)~the lock file lists install/setup steps for the generated skill;
(2)~a subsequent workflow run records the skill as configured, installed, and
available;
(3)~when the task warrants it, the skill is invoked and followed.

Alternatively, the generator may emit \textbf{custom instructions} (prompt-only)
when the repository cannot yet host a skill package; those instructions should
still be structured so a later compile can promote them to \texttt{skills:}
without rewriting policy.

\subsection{Evaluation loop}
We compare three conditions on held-out or replayable triage issues:
\begin{itemize}
  \item inline-prompt triage (no skill);
  \item existing hand-written triage skills (where present);
  \item generated skill compiled via \texttt{gh aw compile}.
\end{itemize}
Primary metrics are label accuracy, override rate, and pickup stages.
Where a full GitHub replay is infeasible, we evaluate generated
\texttt{SKILL.md} intrinsically (SkillEval-style~\cite{skilleval2026}) and
compile-validity (RQ6) as a necessary gate.

%==============================================================================
\section{Study Design}
\label{sec:study-design}

\paragraph{Cohort.}
We combine (i)~periodic gh-aw workflow crawls, (ii)~GitSkills records for
packages referenced by triage workflows, and (iii)~workflow-run artifacts for
repositories that declare triage skills.

\paragraph{Analysis plan.}
RQ1 reports prevalence of triage workflows and skill attachment (repository vs.\
workflow granularity).
RQ2--RQ3 correlate skill features with triage outcomes (regression / ranking).
RQ4 codes approval mechanisms from frontmatter and prompts
(\texttt{issue-intent}, auto-apply, comment reviewer, dry-run).
RQ5--RQ6 implement and evaluate the generator and compile pipeline.

% TODO: Sample size targets, statistical tests, effect-size reporting.

%==============================================================================
\section{Expected Results and Discussion}
\label{sec:results}

% TODO: Replace with empirical findings.

Pilot mining of public gh-aw markdown suggests issue triage is a large share of
agentic workflows, yet skill attachment is rare: most triage policy lives in the
workflow body.
Where skills exist, they encode taxonomies, duplicate procedures, and
suggest-versus-apply rules that inline prompts approximate poorly.
We expect generated skills that include an explicit allowlist, duplicate/spam
sections, and matching \texttt{safe-outputs} to improve label accuracy and to
compile into lock files that actually install the skill---closing the gap
between ``prompt mentions a skill path'' and ``the compiler activated it.''

We also expect approval design (RQ4) to moderate outcomes: suggest-only
\texttt{issue-intent} workflows will show higher override rates but fewer
harmful auto-closes; auto-apply workflows will show faster time-to-triage and
higher need for taxonomy completeness in the skill.

%==============================================================================
\section{Threats to Validity}
\label{sec:threats}

\paragraph{Construct validity.}
``Good'' for triage spans intrinsic skill quality and downstream GitHub utility.
Skill \emph{availability} and \emph{use} are different constructs; conflating
them overestimates effectiveness~\cite{copilotAboutSkills2026}.
Label accuracy requires maintainer ground truth that may lag or disagree.

\paragraph{External validity.}
Public gh-aw repositories may not represent private deployments; generated
skills may overfit mined templates (e.g., \texttt{cli/cli} suggest-only vs.\
SkiaSharp auto-apply).

\paragraph{Internal validity.}
Confounders include repository size, label set complexity, model/engine version,
and whether humans already triaged the issue.

\paragraph{Selection bias.}
Skills declared in frontmatter are rare; early adopters may not generalize.
Keyword-based triage classification can include labeler/PR-triage workflows.

%==============================================================================
\section{Conclusion}
\label{sec:conclusion}

We outlined a mining-and-generation study of \textbf{issue-triage Agent Skills}
in GitHub Agentic Workflows.
The scientific goal is to explain what makes those skills effective; the
engineering goal is a generator (or custom-instruction pack) whose output is
declared so that \texttt{gh aw compile} produces a \filepath{.lock.yml} that
installs and activates the skill in CI.
By linking intrinsic skill properties and sibling artifacts to triage outcomes,
and by treating compile-time activation as a first-class success criterion, we
aim to turn mined maintainer practice into reusable, compiler-ready triage
skills.

%==============================================================================
\balance
\bibliographystyle{ACM-Reference-Format}
\bibliography{references}

\end{document}
